\documentclass[11pt]{article}
\usepackage[T1]{fontenc}
\usepackage[utf8]{inputenc}
\usepackage{lmodern}
\usepackage{microtype}
\usepackage[margin=1in]{geometry}
\usepackage{amsmath,amssymb,amsthm,mathtools}
\usepackage{booktabs,longtable,array}
\usepackage{enumitem}
\usepackage{xcolor}
\usepackage{hyperref}
\usepackage[nameinlink,capitalise,noabbrev]{cleveref}
\usepackage{listings}
\usepackage{fancyhdr}
\usepackage{tikz}
\usetikzlibrary{arrows.meta,positioning,calc}

\hypersetup{
  colorlinks=true,
  linkcolor=blue!55!black,
  citecolor=green!40!black,
  urlcolor=blue!60!black,
  pdftitle={A saturated-lattice and tensor-period reduction for the Alaoglu--Erdos conjecture}
}
\pagestyle{fancy}
\fancyhf{}
\lhead{Alaoglu--Erd\H{o}s: new structural reductions}
\rhead{\thepage}
\setlength{\headheight}{14pt}
\setcounter{tocdepth}{2}
\setcounter{secnumdepth}{3}

\newtheorem{theorem}{Theorem}[section]
\newtheorem{proposition}[theorem]{Proposition}
\newtheorem{lemma}[theorem]{Lemma}
\newtheorem{corollary}[theorem]{Corollary}
\newtheorem{conjecture}[theorem]{Conjecture}
\newtheorem{problem}[theorem]{Problem}
\newtheorem{criterion}[theorem]{Criterion}
\theoremstyle{definition}
\newtheorem{definition}[theorem]{Definition}
\newtheorem{remark}[theorem]{Remark}
\newtheorem{example}[theorem]{Example}

\newcommand{\Q}{\mathbb{Q}}
\newcommand{\R}{\mathbb{R}}
\newcommand{\Z}{\mathbb{Z}}
\newcommand{\N}{\mathbb{N}}
\newcommand{\Gm}{\mathbb{G}_{m}}
\newcommand{\ord}{\operatorname{ord}}
\newcommand{\rank}{\operatorname{rank}}
\newcommand{\Span}{\operatorname{span}}
\newcommand{\divv}{\operatorname{div}}
\newcommand{\cont}{\operatorname{cont}}
\newcommand{\Sym}{\operatorname{Sym}}
\newcommand{\per}{\operatorname{per}}
\newcommand{\supp}{\operatorname{supp}}
\newcommand{\trdeg}{\operatorname{trdeg}}
\newcommand{\Eint}{\mathcal E_{\!\Z}}
\newcommand{\Erat}{\mathcal E_{\!\Q}}
\newcommand{\calD}{\mathcal D}
\newcommand{\calL}{\mathcal L}
\newcommand{\tensorstar}{\mathbin{\odot}}

\title{\textbf{A Saturated-Lattice and Tensor-Period Reduction\\
for the Alaoglu--Erd\H{o}s Conjecture}\\[0.5em]
\large A new progress report centered on exponent semigroups, primitive Pl\"ucker data, and external-prime period orbits}
\author{Prepared for Vladimir Reshetnikov}
\date{23 August 2026}

\begin{document}
\maketitle

\begin{abstract}
Let
\[
  \Eint=\{t\in\R:2^t,3^t\in\N\}.
\]
The Alaoglu--Erd\H{o}s conjecture asserts that \(\Eint=\N_0\).  This report does not repeat the companion report's extensive auxiliary-function, determinant, height, and Diophantine-approximation investigations.  Instead it develops a different structural reduction.

The main unconditional result proved here is a sharp dichotomy.  The six exponentials theorem forces the rational rank of all exponents producing rational values to be at most two.  If a nonintegral exponent exists, saturation of the joint prime-divisor lattice produces a canonical irrational generator \(\xi\), and in fact
\[
  \Eint=\N_0+\N_0\xi
\]
with unique representation.  Thus one counterexample cannot be sporadic: it forces a free rank-two affine semigroup of integer points on the curve \(Y=X^{\log 3/\log 2}\).  This gives an exact counting gap: the conjectural case has \(\asymp\log H\) integer points of height at most \(H\), whereas any counterexample gives \(\asymp(\log H)^2\).

The same saturation is encoded by a primitive quadratic ``two-star'' tensor
\[
 \Sigma=e_2\tensorstar\divv(3^\xi)-e_3\tensorstar\divv(2^\xi)
 \in \Sym^2\!\left(\bigoplus_p\Q e_p\right),
\]
whose real period is zero.  Formal vanishing of \(\Sigma\) gives only the control family \((2^m,3^m)\); a genuine counterexample therefore creates a nonzero primitive kernel tensor with an indispensable external prime.  At such a prime \(q\), the infinitesimal Kummer translation differentiates the alleged quadratic relation to
\(b_q\log2-a_q\log3\ne0\).  This isolates a precise missing theorem: a very small degree-two period-orbit or realization-transfer statement, far narrower than a wholesale proof of Schanuel's conjecture.

No unconditional proof of the Alaoglu--Erd\H{o}s conjecture is claimed.  The report supplies rigorous new reductions, a canonical normal form, a rank-three/rank-four split, an exact polylogarithmic counting target, a concrete period-Galois obstruction, and a formalization blueprint.
\end{abstract}

\tableofcontents
\newpage

\section{Scope, status, and the change of viewpoint}

\subsection{The problem}

\begin{conjecture}[Alaoglu--Erd\H{o}s]
For real \(x\),
\[
  2^x\in\Z,\qquad 3^x\in\Z
  \quad\Longrightarrow\quad x\in\Z.
\]
Since positive real powers are positive, the integers in the hypothesis are automatically positive.
\end{conjecture}

Put
\[
 A=2^x,\qquad B=3^x,
 \qquad \alpha=\frac{\log3}{\log2}.
\]
Then \(B=A^\alpha\), and the conjecture says that the integer points on
\[
  C_\alpha:\quad Y=X^\alpha,\qquad X,Y\ge1,
\]
are exactly \((2^m,3^m)\), \(m\in\N_0\).

The familiar determinant is
\begin{equation}\label{eq:basic-det}
 \det\begin{pmatrix}\log2&\log A\\[2pt]\log3&\log B\end{pmatrix}
 =\log2\log B-\log3\log A=0.
\end{equation}
For nonintegral \(x\), both the row and column rational ranks are two, so \eqref{eq:basic-det} is exactly a rational, positive-integer instance of the four exponentials problem.  Rephrasing \eqref{eq:basic-det} alone therefore does not constitute progress.

The companion report has already explored many quantitative incarnations of this determinant.  The present report changes the object under study: instead of asking first how small a determinant can be, it asks what the \emph{whole set of admissible exponents} and its prime-divisor lattice would have to look like.  This produces a rigid global normal form before any transcendence estimate is attempted.

\subsection{What is proved here}

The genuinely new structural chain is:
\begin{enumerate}[label=(\roman*),leftmargin=2.4em]
\item the six exponentials theorem bounds the rational rank of all admissible exponents by two;
\item a valuation embedding turns admissible exponents into a submonoid of a free abelian divisor lattice;
\item saturation extracts a canonical primitive irrational generator \(\xi\);
\item positivity at an external prime and at one distinguished base place forces every admissible exponent to be \(m+n\xi\) with \(m,n\ge0\);
\item the same saturation index is the content of a sparse quadratic period tensor;
\item an external prime differentiates that tensor to a nonzero linear logarithmic form.
\end{enumerate}
Items (i)--(v) are unconditional.  Item (vi) supplies a sharp obstruction and a conditional closing lemma, but the required realization-transfer theorem is not currently established here.

\subsection{Baseline reductions}

We record the elementary facts needed later.

\begin{lemma}\label{lem:baseline}
If \(x\in\Eint\), then:
\begin{enumerate}[label=(\alph*)]
\item \(x\ge0\);
\item if \(x\in\Q\), then \(x\in\N_0\);
\item if \(x\notin\Z\), then \(x\) is transcendental;
\item \(A\) and \(B\) are multiplicatively independent.
\end{enumerate}
\end{lemma}

\begin{proof}
If \(x<0\), then \(0<2^x<1\), proving (a).  Write a rational \(x=r/s\) in lowest terms with \(s>0\).  If \(2^{r/s}=A\in\N\), then \(A^s=2^r\); comparison of the \(2\)-adic valuations gives \(s\mid r\), hence \(s=1\), proving (b).  If a nonintegral \(x\) were algebraic, it would be algebraic irrational, and the Gelfond--Schneider theorem would make \(2^x\) transcendental, proving (c).  Finally, \(A^u=B^v\) with nonzero integers \(u,v\) would give \(2^u=3^v\) after taking the positive \(x\)-th root, impossible by unique factorization.  This proves (d).
\end{proof}

The number \(\alpha\) is itself transcendental: it is irrational by unique factorization, while an algebraic irrational \(\alpha\) would make \(2^\alpha=3\) transcendental by Gelfond--Schneider.

\section{The divisor-lattice model}

\subsection{Formal logarithms}

Let
\[
 \calD=\bigoplus_{p\ \mathrm{prime}}\Z e_p,
 \qquad
 \calD_\Q=\calD\otimes_\Z\Q.
\]
For a positive rational number \(r\), write
\[
 \divv(r)=\sum_p v_p(r)e_p\in\calD.
\]
The real logarithm is the \(\Q\)-linear map
\[
 \ell:\calD_\Q\longrightarrow\R,
 \qquad \ell(e_p)=\log p.
\]
Unique factorization says that \(\ell\) is injective: a rational linear relation among the \(\log p\)'s can be cleared of denominators and exponentiated to a multiplicative relation among distinct primes.

For pairs of rational numbers use
\[
 F=\calD\oplus\calD,
 \qquad
 \nu(t)=\bigl(\divv(2^t),\divv(3^t)\bigr)

to be defined whenever both powers are positive rationals.

\begin{definition}
Define the rational-value exponent group and integer-value exponent monoid by
\[
 \Erat=\{t\in\R:2^t,3^t\in\Q_{>0}\},
 \qquad
 \Eint=\{t\in\R:2^t,3^t\in\N\}.
\]
\end{definition}

The map \(\nu:\Erat\to F\) is an injective group homomorphism.  It is additive because divisors are additive under multiplication, and it is injective because \(2^t=1\) implies \(t=0\).  Moreover,
\[
 \nu(\Eint)=\nu(\Erat)\cap(F_{\ge0}),
\]
where \(F_{\ge0}\) denotes coordinatewise nonnegative divisor vectors.

\subsection{The two distinguished vectors}

Let
\[
 g_0=\nu(1)=(e_2,e_3).
\]
For a putative counterexample \(x\), write
\[
 a=\divv(A)=\sum_pa_pe_p,
 \qquad
 b=\divv(B)=\sum_pb_pe_p,
 \qquad
 g_x=(a,b).
\]
The real-log map sends these vectors to
\[
 (\ell\oplus\ell)(g_0)=(\log2,\log3),
 \qquad
 (\ell\oplus\ell)(g_x)=x(\log2,\log3).
\]
The formal vectors \(g_0,g_x\) are nevertheless \(\Q\)-linearly independent when \(x\notin\Q\): applying real logs to a rational relation would give the same rational relation between \(1\) and \(x\).

\section{Six exponentials gives a global rank cap}

\begin{theorem}[Six exponentials theorem]\label{thm:six-exp}
Let \(u_1,u_2\in\C\) be linearly independent over \(\Q\), and let \(v_1,v_2,v_3\in\C\) be linearly independent over \(\Q\).  Then at least one of the six numbers
\[
 e^{u_i v_j}\qquad(1\le i\le2,\ 1\le j\le3)
\]
is transcendental.
\end{theorem}

\begin{proposition}[Rank cap]\label{prop:rank-cap}
The \(\Q\)-vector space spanned by \(\Erat\) has dimension at most two.
\end{proposition}

\begin{proof}
The numbers \(\log2\) and \(\log3\) are \(\Q\)-linearly independent.  If \(t_1,t_2,t_3\in\Erat\) were \(\Q\)-linearly independent, apply \cref{thm:six-exp} with
\[
 (u_1,u_2)=(\log2,\log3),
 \qquad (v_1,v_2,v_3)=(t_1,t_2,t_3).
\]
All six exponentials would be positive rational numbers, a contradiction.
\end{proof}

\begin{corollary}\label{cor:rank-dichotomy}
The group \(\Erat\), embedded by \(\nu\) in the free abelian group \(F\), is a free abelian group of rank one or two.  It contains \(\Z\), and \(1\) is primitive in \(\Erat\).
\end{corollary}

\begin{proof}
A subgroup of a free abelian group is free.  Its rank is the dimension of its rational span, bounded by \cref{prop:rank-cap}.  It has rank at least one because it contains \(1\).  If \(1=d t\) in \(\Erat\) for an integer \(d>1\), then \(2^t\) would be a rational \(d\)-th root of \(2\), impossible by valuations.  Thus \(1\) is primitive.
\end{proof}

This is the first important global restriction.  Four exponentials asks whether rank two is possible; six exponentials says rank three is already impossible.  Consequently, once one nonrational exponent is fixed, every other rational-value exponent lies in its rational plane.

\section{Saturation and the canonical primitive counterexample}

\subsection{The saturated plane}

Fix \(x\in\Eint\setminus\Q\).  Define
\[
 W_x=\Q g_0+\Q g_x\subset F_\Q,
 \qquad
 \Lambda_x=W_x\cap F.
\]
The lattice \(\Lambda_x\) is the saturation of \(\Z g_0+\Z g_x\) in \(F\).  Since \(g_0\) is primitive, it can be completed to a basis
\[
 \Lambda_x=\Z g_0\oplus\Z h.
\]
There are unique integers \(r\) and \(d>0\), after fixing the sign of \(h\), such that
\begin{equation}\label{eq:smith}
 g_x=r g_0+d h.
\end{equation}
Replacing \(h\) by \(h+k g_0\) replaces \(r\) by \(r-dk\).  This harmless operation is the divisor-lattice version of shifting an exponent by an integer.

\begin{proposition}[Canonical saturated generator]\label{prop:canonical-generator}
If \(\Eint\ne\N_0\), there exists an irrational \(\xi\in\Eint\) such that, with
\[
 C=2^\xi,\qquad D=3^\xi,
\]
the following hold.
\begin{enumerate}[label=(\roman*)]
\item \(\Lambda=\bigl(\Q\nu(1)+\Q\nu(\xi)\bigr)\cap F
      =\Z\nu(1)\oplus\Z\nu(\xi)\).
\item \(\min\{v_2(C),v_3(D)\}=0\).
\item At least one prime \(q\notin\{2,3\}\) divides \(CD\).
\end{enumerate}
\end{proposition}

\begin{proof}
Start with a nonrational \(x\in\Eint\) and the basis relation \eqref{eq:smith}.  At every coordinate other than the first-copy \(2\)-coordinate and the second-copy \(3\)-coordinate, \(g_0\) is zero.  Hence those coordinates of \(h\) are the corresponding nonnegative coordinates of \(g_x\), divided by \(d\).  By replacing \(h\) with \(h+k g_0\) for sufficiently large \(k\), its two distinguished coordinates also become nonnegative.  Then choose \(k\) uniquely so that the smaller of those two coordinates is zero.  All coordinates of the resulting \(h\) are nonnegative.

Because \(h\in W_x\), write \(h=u g_0+v g_x\) with \(u,v\in\Q\) and set \(\xi=u+vx\).  Applying the injective formal-log map in each coordinate shows that the two divisor vectors forming \(h\) are exactly \(\divv(2^\xi)\) and \(\divv(3^\xi)\).  Nonnegativity makes these powers integers.  The basis property gives (i), and the chosen shift gives (ii).  The number \(\xi\) is irrational, since \(x=r+d\xi\) and a rational \(\xi\) would make \(x\) rational.

Suppose (iii) failed.  Write
\[
 C=2^{a_2}3^{a_3},\qquad D=2^{b_2}3^{b_3}.
\]
Dividing the two logarithmic identities by \(\log2\) gives
\[
 \xi=a_2+a_3\alpha,
 \qquad
 \xi\alpha=b_2+b_3\alpha.
\]
Thus
\[
 a_3\alpha^2+(a_2-b_3)\alpha-b_2=0.
\]
The transcendence of \(\alpha\) forces \(a_3=b_2=0\) and \(a_2=b_3\), whence \(\xi=a_2\in\Z\), a contradiction.
\end{proof}

The generator \(\xi\) will be called \emph{canonical}.  Property (i) is the saturation condition, property (ii) fixes the integer-shift gauge, and property (iii) supplies the external-prime witness.

\subsection{An explicit content formula}

For arbitrary \(A=\prod p^{a_p}\), \(B=\prod p^{b_p}\), define
\begin{equation}\label{eq:content}
 \delta(A,B)=\gcd\Bigl(
   \{a_p:p\ne2\}\cup
   \{b_p:p\ne3\}\cup
   \{b_3-a_2\}
 \Bigr),
\end{equation}
where the gcd is nonnegative and zeros are harmless.

\begin{proposition}[Saturation index]\label{prop:content-index}
For \(g_0=(e_2,e_3)\) and \(g_x=(a,b)\), the index
\[
 [\,(\Q g_0+\Q g_x)\cap F:\Z g_0+\Z g_x\,]
\]
is \(\delta(A,B)\).  In particular, the pair is saturated exactly when \(\delta(A,B)=1\).
\end{proposition}

\begin{proof}
For two integer vectors, the index of their span in its saturation is the gcd of all \(2\times2\) minors of the matrix having those vectors as columns.  Since \(g_0\) has only two nonzero entries, every nonzero minor is, up to sign, one of
\[
 a_p\ (p\ne2),\qquad b_p\ (p\ne3),\qquad b_3-a_2.
\]
Their gcd is \eqref{eq:content}.
\end{proof}

\begin{corollary}[Explicit root descent]\label{cor:root-descent}
If \(x\in\Eint\setminus\Z\) and \(d=\delta(A,B)>1\), choose an integer \(r\) satisfying
\[
 r\equiv a_2\equiv b_3\pmod d,
 \qquad r\le\min(a_2,b_3).
\]
Then
\[
 A_1=\left(\frac{A}{2^r}\right)^{1/d},
 \qquad
 B_1=\left(\frac{B}{3^r}\right)^{1/d}
\]
are positive integers and
\[
 x_1=\frac{x-r}{d}\in\Eint\setminus\Z.
\]
The divisor pair for \(x_1\) is saturated.
\end{corollary}

\begin{proof}
The congruences and the definition of \(d\) make every prime exponent in \(A/2^r\) and \(B/3^r\) divisible by \(d\); the inequality on \(r\) makes the distinguished exponents nonnegative.  The displayed power identities follow from positive real roots.  Dividing all Pl\"ucker coordinates by their gcd makes the new lattice primitive.
\end{proof}

This descent is not a size descent on \(A\) or \(B\); it is a canonical descent on the index of the exponent lattice.  It is therefore immune to the absence of a least transcendental counterexample.

\section{The exponent-semigroup dichotomy}

We now prove the strongest unconditional structural result of this report.

\begin{theorem}[Free-semigroup dichotomy]\label{thm:semigroup-dichotomy}
Exactly one of the following holds.
\begin{enumerate}[label=(\mathrm{I}),leftmargin=3em]
\item \(\Eint=\N_0\).
\item There is a canonical irrational \(\xi>1\), unique with the normalization of \cref{prop:canonical-generator}, such that
\begin{equation}\label{eq:E-free}
 \Eint=\N_0+\N_0\xi.
\end{equation}
Every element has a unique representation \(m+n\xi\), \(m,n\in\N_0\).
\end{enumerate}
\end{theorem}

\begin{proof}
Assume the conjecture is false and choose the canonical \(\xi\) of \cref{prop:canonical-generator}.  The inclusion from right to left in \eqref{eq:E-free} is immediate:
\[
 2^{m+n\xi}=2^m(2^\xi)^n\in\N,
 \qquad
 3^{m+n\xi}=3^m(3^\xi)^n\in\N.
\]

Conversely, let \(y\in\Eint\).  By \cref{prop:rank-cap}, \(y\in\Q+\Q\xi\); write \(y=u+v\xi\).  Let \(h_y=\nu(y)\).  In the first divisor coordinate,
\[
 \ell\bigl(\divv(2^y)-u e_2-v\divv(2^\xi)\bigr)=0.
\]
The vector inside \(\ell\) has rational coefficients, so injectivity of \(\ell\) makes it zero.  The same argument applies to the second coordinate.  Hence
\[
 h_y=u\nu(1)+v\nu(\xi)\in\Lambda.
\]
The saturated basis property gives integers \(m,n\) with
\[
 h_y=m\nu(1)+n\nu(\xi),
 \qquad y=m+n\xi.
\]

At an external prime coordinate where \(\nu(\xi)\) is positive and \(\nu(1)\) is zero, nonnegativity of \(h_y\) forces \(n\ge0\).  At whichever distinguished coordinate realizes
\(\min\{v_2(2^\xi),v_3(3^\xi)\}=0\), the vector \(\nu(\xi)\) is zero and \(\nu(1)\) is one, so nonnegativity forces \(m\ge0\).  This proves \eqref{eq:E-free}.  Uniqueness of the representation follows from the irrationality of \(\xi\).

Finally \(\xi>1\), because \(0<\xi<1\) would put the integer \(2^\xi\) strictly between \(1\) and \(2\).  In the free monoid \(\N_0+\N_0\xi\), \(\xi\) is the unique nonintegral irreducible element; the normalization therefore makes it unique.
\end{proof}

\begin{corollary}[No sporadic counterexamples]\label{cor:no-sporadic}
If one nonintegral solution exists, then there are quadratically many admissible exponents up to size \(T\), and every admissible exponent is generated additively by \(1\) and a single canonical counterexample.
\end{corollary}

\begin{corollary}[Exact lattice-point asymptotic]\label{cor:count-T}
In case (II) of \cref{thm:semigroup-dichotomy},
\[
 \#\{t\in\Eint:0\le t\le T\}
 =\frac{T^2}{2\xi}+O_\xi(T).
\]
In case (I), the count is \(\lfloor T\rfloor+1\).
\end{corollary}

\begin{proof}
In case (II), sum over \(0\le n\le T/\xi\):
\[
 \sum_{0\le n\le T/\xi}
 \bigl(\lfloor T-n\xi\rfloor+1\bigr)
 =\frac{T^2}{2\xi}+O_\xi(T).
\]
The other case is immediate.
\end{proof}

This theorem sharply reorganizes the problem.  A hypothetical universe with one extra solution is not a universe with a rare accident; it has a normal affine semigroup of rank two sitting on a one-dimensional transcendental curve.

\section{The primitive two-star tensor}

\subsection{Quadratic formal periods}

Pass to the rational prime-divisor space
\[
 V=\calD_\Q=\bigoplus_p\Q e_p.
\]
Let \(\Sym^2(V)\) be the commutative tensor square.  We write \(u\tensorstar v\) for the commutative monomial represented by \(uv\), with no factor of \(1/2\).  Define the quadratic real-period map
\[
 \per_2:\Sym^2(V)\longrightarrow\R,
 \qquad
 \per_2(e_p\tensorstar e_q)=\log p\log q.
\]

For \(A=\prod p^{a_p}\), \(B=\prod p^{b_p}\), define
\begin{equation}\label{eq:star-tensor}
 \Sigma(A,B)=e_2\tensorstar b-e_3\tensorstar a.
\end{equation}
Its period is
\begin{equation}\label{eq:star-period}
 \per_2(\Sigma(A,B))=\log2\log B-\log3\log A.
\end{equation}
Thus an admissible exponent produces a kernel tensor.

Expanding in the prime basis gives the sparse ``two-star'' form
\begin{equation}\label{eq:star-expanded}
\begin{split}
 \Sigma(A,B)={}&b_2 e_2^2-a_3e_3^2+(b_3-a_2)e_2e_3\\
 &+\sum_{p\notin\{2,3\}}
   \bigl(b_p e_2e_p-a_p e_3e_p\bigr).
\end{split}
\end{equation}
Its support graph consists of a star centered at \(2\) and a star centered at \(3\), with opposite signs on the two families of external edges.

\begin{proposition}[Formal zeroes are exactly the controls]\label{prop:formal-zero}
One has \(\Sigma(A,B)=0\) in \(\Sym^2(V)\) if and only if
\[
 A=2^m,\qquad B=3^m
\]
for some \(m\in\N_0\).
\end{proposition}

\begin{proof}
Comparing the coefficients in \eqref{eq:star-expanded}, formal vanishing gives
\[
 b_2=0,\quad a_3=0,\quad b_3=a_2,
 \quad a_p=b_p=0\ (p\notin\{2,3\}).
\]
This is precisely the stated control family.  The converse follows from commutativity: \(e_2e_3=e_3e_2\).
\end{proof}

The crucial distinction is now visible.  For the control solution \(x=m\), the numerical determinant vanishes because its formal symmetric tensor already vanishes.  For a counterexample, the same numerical vanishing would kill a \emph{nonzero} formal tensor.

\subsection{Content equals saturation}

For an integral tensor, let its content be the gcd of its prime-basis coefficients.  From \eqref{eq:star-expanded},
\[
 \cont\Sigma(A,B)
 =\gcd\Bigl(
 b_2,a_3,b_3-a_2,
 \{a_p,b_p:p\notin\{2,3\}\}
 \Bigr),
\]
which is the same number as \(\delta(A,B)\) in \eqref{eq:content}.

\begin{theorem}[Primitive tensor normal form]\label{thm:tensor-normal}
If the Alaoglu--Erd\H{o}s conjecture is false, then there are nonnegative finite exponent vectors \(a,b\) such that:
\begin{enumerate}[label=(\roman*)]
\item \(\per_2(e_2\tensorstar b-e_3\tensorstar a)=0\);
\item \(e_2\tensorstar b-e_3\tensorstar a\ne0\);
\item its content is one;
\item \(\min(a_2,b_3)=0\);
\item some external coefficient \(a_q\) or \(b_q\), \(q\notin\{2,3\}\), is positive.
\end{enumerate}
Conversely, any such pair \((a,b)\) reconstructs a nonintegral counterexample by
\[
 \xi=\frac{\ell(a)}{\log2}=\frac{\ell(b)}{\log3}.
\]
\end{theorem}

\begin{proof}
Use the canonical \(\xi\) and its divisor vectors.  The period vanishing is \eqref{eq:star-period}; nonzero follows from \cref{prop:formal-zero} and irrationality; primitivity follows from saturation and \cref{prop:content-index}; the last two properties are part of \cref{prop:canonical-generator}.  Conversely, the period identity makes the two displayed ratios equal.  The nonnegative exponent vectors make the powers integers.  If the common ratio were integral, formal uniqueness of prime factorization would give the control tensor zero, contrary to (ii).
\end{proof}

\begin{criterion}[Star-tensor criterion]\label{crit:star}
The Alaoglu--Erd\H{o}s conjecture is equivalent to the assertion that the kernel of \(\per_2\) contains no primitive nonzero integral two-star tensor satisfying the sign and normalization conditions of \cref{thm:tensor-normal}.
\end{criterion}

This criterion is not merely a renaming of \eqref{eq:basic-det}.  It removes three symmetries that otherwise pollute auxiliary constructions:
\begin{itemize}[leftmargin=2em]
\item the commutativity relation responsible for the control family;
\item integer shifts \(x\mapsto x+k\);
\item simultaneous perfect-power extraction.
\end{itemize}
What remains is a primitive, externally supported quadratic period relation.

\section{Further unconditional structure}

\subsection{Multiplicative rank is at least three}

Let
\[
 \rho=\dim_\Q\Span_\Q\{e_2,e_3,a,b\}.
\]
This is the multiplicative rank of \(2,3,A,B\).

\begin{proposition}\label{prop:rank-at-least-3}
For a nonintegral solution, \(\rho\in\{3,4\}\).
\end{proposition}

\begin{proof}
The rank is at most four.  If \(\rho\le2\), then \(a,b\in\Q e_2+\Q e_3\).  Since they are integral nonnegative vectors, write
\[
 A=2^r3^s,\qquad B=2^u3^v
\]
with nonnegative integers \(r,s,u,v\).  The two power identities give
\[
 x=r+s\alpha,
 \qquad x\alpha=u+v\alpha,
\]
so
\[
 s\alpha^2+(r-v)\alpha-u=0.
\]
Transcendence of \(\alpha\) forces \(s=u=0\), \(r=v\), and \(x=r\in\Z\), a contradiction.
\end{proof}

\subsection{The rank-three normal form}

Let \(\bar a,\bar b\) be the images of \(a,b\) in
\(V/(\Q e_2+\Q e_3)\).  Rank three means that these two external divisor vectors span a line.

\begin{proposition}[Rank-three normal form]\label{prop:rank3-normal}
If \(\rho=3\), there exist a positive integer \(C\) coprime to \(6\), not a proper perfect power, nonnegative integers
\[
 a,c,d,b,u,v,
 \qquad (u,v)\ne(0,0),
\]
such that
\begin{equation}\label{eq:rank3-factor}
 A=2^a3^cC^u,
 \qquad
 B=2^d3^bC^v.
\end{equation}
Writing \(\gamma=\log C/\log2\), one has
\[
 x=a+c\alpha+u\gamma,
 \qquad
 x\alpha=d+b\alpha+v\gamma,
\]
and hence
\begin{equation}\label{eq:mobius}
 x=\frac{(va-ud)+(vc-ub)\alpha}{v-u\alpha}.
\end{equation}
The denominator is nonzero.
\end{proposition}

\begin{proof}
Choose the primitive nonnegative generator of the one-dimensional lattice spanned by the external projections of \(a,b\); it is the divisor of a unique integer \(C\) coprime to \(6\) and not a proper perfect power.  This gives \eqref{eq:rank3-factor}.  The logarithmic equations are immediate.  Multiplying the first by \(v\), the second by \(u\), and subtracting eliminates \(\gamma\), yielding \eqref{eq:mobius}.  If \(v-u\alpha=0\), irrationality of \(\alpha\) forces \(u=v=0\), excluded.
\end{proof}

Thus every rank-three counterexample lies in a one-transcendence-parameter normal form: \(x\) is a rational fractional-linear function of \(\alpha\).  This is a sharply delimited subproblem.  It is not currently eliminated by Gelfond--Schneider; a fractional-linear function of a transcendental number may remain transcendental, and the algebraicity of the four exponentials is exactly the hard input.

Rank four is the complementary case in which the external parts of \(A\) and \(B\) are multiplicatively independent modulo powers of \(2\) and \(3\).  Any proposed argument should state explicitly which of these two cases it handles.  Many reductions that appear to gain an extra independent prime silently fail in rank three, while many one-parameter algebraic-independence ideas do not address rank four.

\subsection{Slope straddling}

The determinant relation can also be written
\begin{equation}\label{eq:slope-balance}
 \sum_p (b_p-a_p\alpha)\log p=0.
\end{equation}

\begin{lemma}[Slope straddling]\label{lem:straddle}
For a nonintegral solution, the nonzero coefficients \(b_p-a_p\alpha\) in \eqref{eq:slope-balance} cannot all have the same sign.  Equivalently, the local valuation slopes \(b_p/a_p\), with the usual conventions at zero, occur on both sides of \(\alpha\).
\end{lemma}

\begin{proof}
Every \(\log p\) is positive.  A nontrivial weighted sum with all coefficients of one weak sign cannot vanish.  No coefficient can vanish with \((a_p,b_p)\ne(0,0)\), because that would make \(\alpha=b_p/a_p\) rational.
\end{proof}

This elementary oriented-matroid constraint is useful in computations and in choosing an external prime at which a period-orbit derivative has a prescribed sign.

\section{A sharp counting reformulation}

\subsection{Integer points on the irrational power curve}

Let
\[
 N_\alpha(H)=\#\{(M,N)\in\N^2:N=M^\alpha,\ \max(M,N)\le H\}.
\]
Every point has the form
\[
 M=2^t,\qquad N=3^t,
 \qquad t=\frac{\log M}{\log2}\in\Eint.
\]
Since \(\alpha>1\), the height condition is equivalent to
\(t\le(\log H)/(\log3)\).

\begin{theorem}[Polylogarithmic growth gap]\label{thm:growth-gap}
If the Alaoglu--Erd\H{o}s conjecture is true, then
\[
 N_\alpha(H)=\frac{\log H}{\log3}+O(1).
\]
If it is false and \(\xi\) is the canonical generator, then
\begin{equation}\label{eq:height-asymptotic}
 N_\alpha(H)=
 \frac{(\log H)^2}{2\xi(\log3)^2}+O_\xi(\log H).
\end{equation}
\end{theorem}

\begin{proof}
Combine \cref{thm:semigroup-dichotomy,cor:count-T} with
\(T=(\log H)/(\log3)\).
\end{proof}

\begin{corollary}[A counting target]\label{cor:counting-target}
Any unconditional estimate
\[
 N_\alpha(H)=o\bigl((\log H)^2\bigr)
\]
proves the Alaoglu--Erd\H{o}s conjecture.
\end{corollary}

This target is qualitatively different from proving that a single determinant is nonzero.  It asks for a small improvement over quadratic polylogarithmic growth.  Standard Pila--Wilkie counting for definable transcendental curves gives \(O_\varepsilon(H^\varepsilon)\), which is vastly larger than every fixed power of \(\log H\) and therefore does not reach \cref{cor:counting-target}.  A sharp Wilkie-type estimate with exponent strictly below two for this particular exponential curve would suffice.

The dichotomy also says that an upper bound merely excluding \emph{rank-two} polylogarithmic growth is enough.  One need not classify all individual points.  This may fit determinant methods that count many points more naturally than methods designed to certify one fixed logarithmic determinant.

\subsection{Why the lower bound is especially rigid}

A false conjecture gives not just arbitrary points but the explicit grid
\[
 (M_{m,n},N_{m,n})=(2^mC^n,3^mD^n),
 \qquad m,n\in\N_0,
\]
where \(C=2^\xi\), \(D=3^\xi\).  Distinct pairs \((m,n)\) give distinct points because \(\xi\) is irrational.  Their logarithmic coordinates form a unimodular lattice cone:
\[
 (\log M_{m,n},\log N_{m,n})
 =(m+n\xi)(\log2,\log3).
\]
Any counting proof can therefore exploit more than cardinality: it may use the multiplicative recurrences
\[
 M_{m+1,n}=2M_{m,n},\quad M_{m,n+1}=CM_{m,n},
\]
\[
 N_{m+1,n}=3N_{m,n},\quad N_{m,n+1}=DN_{m,n}.
\]
A promising determinant should be built on this two-dimensional recurrence grid rather than on unrelated sampled points of the curve.

\section{The external-prime period-orbit obstruction}

\subsection{A formal Kummer translation}

Fix a finite prime set \(S\) containing the supports of \(6AB\).  Introduce formal variables \(X_p\) for \(p\in S\) and a period symbol \(\Omega\), heuristically specialized to
\[
 X_p=\log p,\qquad \Omega=2\pi i.
\]
Set
\[
 L_A=\sum_pa_pX_p,
 \qquad
 L_B=\sum_pb_pX_p,
 \qquad
 F=X_2L_B-X_3L_A.
\]
The counterexample relation is \(F(\log p)=0\).

For an external prime \(q\notin\{2,3\}\), consider the infinitesimal translation
\[
 X_q\longmapsto X_q+t\Omega,
 \qquad X_p\longmapsto X_p\quad(p\ne q).
\]
Its derivative acts by \(\Omega\partial_q\).  A direct calculation gives
\begin{equation}\label{eq:q-derivative}
 \partial_qF=b_qX_2-a_qX_3.
\end{equation}
At the real logarithms this is nonzero whenever \((a_q,b_q)\ne(0,0)\), because equality would imply
\(2^{b_q}=3^{a_q}\).

\begin{proposition}[Orbit obstruction]\label{prop:orbit-obstruction}
A primitive counterexample relation is not invariant under the independent Kummer translation at any external prime in its support.  Its first orbit derivative is the nonzero linear logarithmic form
\[
 \Omega\,(b_q\log2-a_q\log3).
\]
\end{proposition}

This one-line derivative is the reason to insist on an external prime.  At \(2\) or \(3\), differentiating also hits the two distinguished factors and mixes in control terms.  At an external \(q\), the derivative cleanly reads one pair of valuation coefficients.

\subsection{The exact conditional closing lemma}

\begin{problem}[Degree-two Kummer orbit principle]\label{prob:kummer-orbit}
Let \(S\) be a finite set of rational primes.  Suppose a quadratic polynomial in \(\Omega\) and the Kummer logarithmic periods \(X_p\) vanishes at
\((2\pi i,\log p)\).  Is every such vanishing relation stable under the independent unipotent Kummer translations
\(X_q\mapsto X_q+t_q\Omega\) dictated by the motivic Galois group?
\end{problem}

\begin{theorem}[Conditional closure]\label{thm:conditional-orbit}
A positive answer to \cref{prob:kummer-orbit} for the two-star quadratics \(F=X_2L_B-X_3L_A\) proves the Alaoglu--Erd\H{o}s conjecture.
\end{theorem}

\begin{proof}
A counterexample gives a primitive \(F\) with an external prime \(q\) in its support.  Stability of the relation ideal under the \(q\)-translation would force its orbit derivative to vanish at the periods.  By \eqref{eq:q-derivative}, this says
\(b_q\log2-a_q\log3=0\), contrary to multiplicative independence of \(2\) and \(3\).
\end{proof}

This principle should not be confused with the known linear theory of logarithms or with period results for linear relations among \(1\)-periods.  The relation \(F=0\) is a product-of-logarithms relation, hence lives in a tensor construction.  Linear period theorems do not automatically transport it through the full unipotent orbit.  The missing statement is precisely a degree-two, externally supported matrix-coefficient assertion.

The gain in formulation is substantial.  One need not establish algebraic independence of all prime logarithms, nor the whole four exponentials conjecture.  It is enough to show that a zero of this single sparse two-star matrix coefficient cannot disappear under an independently detectable Kummer direction.

\section{Finite-place witnesses and Kummer primitivity}

\subsection{Reduction modulo an auxiliary prime}

For a canonical pair \((C,D)=(2^\xi,3^\xi)\), write exponent vectors \(a,b\).  Content one is equivalent to the following statement for every rational prime \(\ell\): there is no residue \(r\in\Z/\ell\Z\) such that
\begin{equation}\label{eq:shifted-power-congruence}
 a\equiv r e_2\pmod\ell,
 \qquad
 b\equiv r e_3\pmod\ell.
\end{equation}
Indeed, \eqref{eq:shifted-power-congruence} says that after the common integer shift \(\xi\mapsto\xi-r\), both integer values are simultaneous \(\ell\)-th powers.

\begin{proposition}[Kummer primitivity]\label{prop:kummer-primitive}
For every prime \(\ell\), the two divisor vectors \(\nu(1)\) and \(\nu(\xi)\) are linearly independent in \(F/\ell F\).  Equivalently, at least one of the following witnesses occurs:
\begin{enumerate}[label=(\alph*)]
\item \(a_p\not\equiv0\pmod\ell\) for some \(p\ne2\);
\item \(b_p\not\equiv0\pmod\ell\) for some \(p\ne3\);
\item \(a_2\not\equiv b_3\pmod\ell\).
\end{enumerate}
\end{proposition}

\begin{proof}
Linear dependence modulo \(\ell\) would make every \(2\times2\) minor of the two-column divisor matrix zero modulo \(\ell\).  Those minors are exactly the coefficients whose gcd is the content of \(\Sigma(C,D)\), which is one.
\end{proof}

This supplies a finite-place witness at every \(\ell\).  Kummer theory can turn such valuation data into nontrivial power-residue characters in suitable extensions.  What is missing is a theorem forcing the archimedean quadratic period relation to respect one of those finite-place witnesses.  The desired transfer would have the schematic form
\[
 \text{complex two-star period }=0
 \quad\Longrightarrow\quad
 b_q\log2-a_q\log3=0
\]
for some externally detected \(q\), which is impossible.

\subsection{Why naive \(p\)-adic substitution does not work}

It is tempting to replace every real logarithm in \eqref{eq:basic-det} by a \(p\)-adic logarithm.  This is invalid.  The real identity
\[
 \log A=x\log2,
 \qquad \log B=x\log3
\]
does not furnish a common \(p\)-adic number \(x\) satisfying the analogous identities.  The number \(x\) is defined by an archimedean quotient of logarithms, and there is no canonical comparison with \(\log_p A/\log_p2\).  A successful mixed-place argument therefore needs a genuine comparison or period-orbit theorem; it cannot simply change absolute values in the displayed determinant.

This observation is a useful filter.  Any proposed local proof should identify the exact functorial object carrying the relation from the complex realization to the finite-place realization.  Without that object, the argument has silently assumed the missing theorem.

\section{A determinant design adapted to the semigroup}

The companion report already analyzes many auxiliary determinants.  The structural results above impose four tests on any new one.

\subsection{Gauge invariance}

Replacing \(x\) by \(x+k\) multiplies \((A,B)\) by \((2^k,3^k)\) and leaves \(\Sigma(A,B)\) unchanged.  A determinant that changes essentially under this operation is measuring a choice of representative rather than the obstruction.

\subsection{Root-descent invariance}

If the Pl\"ucker content is \(d>1\), \cref{cor:root-descent} extracts a smaller saturated exponent.  An auxiliary construction should be homogeneous enough that this extraction divides its obstruction rather than creating artificial denominators.

\subsection{Control annihilation}

For \((A,B)=(2^m,3^m)\), the symmetric tensor is formally zero.  A viable determinant should annihilate this family algebraically before estimates are made.  If the control family only cancels numerically at the final line, the same cancellation is likely to survive for a counterexample.

\subsection{External-prime sensitivity}

At least one row, column, derivative, or local factor should see a coefficient \(a_q\) or \(b_q\) for \(q\notin\{2,3\}\).  The cleanest target is the nonzero derivative \(b_q\log2-a_q\log3\).  A construction depending only on \(\log A=x\log2\) and \(\log B=x\log3\) cannot distinguish a counterexample from the controls.

\subsection{A grid rather than a line}

Under a counterexample, the full set of integer points is
\[
 (2^mC^n,3^mD^n),\qquad m,n\ge0.
\]
One should therefore interpolate on a two-dimensional recurrence grid and use its unimodularity.  A one-dimensional sample \(t=n\xi\) discards the ordinary generator \(1\), while a sample \(t=m\) discards the counterexample direction.  The rank-two grid is the first place where the false and true worlds have different growth.

A possible implementation is an interpolation determinant built from monomials \(X^iY^j\) evaluated on a triangular set of grid points.  The determinant factors into products of
\[
 2^{m_1}C^{n_1}-2^{m_2}C^{n_2}
 \quad\text{and}\quad
 3^{m_1}D^{n_1}-3^{m_2}D^{n_2},
\]
so it is an integer, while the points are constrained to the same analytic curve.  The unresolved quantitative issue is to make the analytic upper bound beat the enormous product of heights without reintroducing the usual four-exponentials barrier.  The counting target of \cref{cor:counting-target} suggests optimizing for a subquadratic number of effective degrees of freedom rather than for a single spectacularly small determinant.

\section{Exploratory computation}

Numerics cannot certify the absence of an exact product-of-logarithms relation in an unbounded range, but they are useful for testing scale choices in an auxiliary construction.  The following experiment scanned
\[
 2\le A\le 1,000,000,
\]
excluding powers of two, and ranked the distance from \(A^\alpha\) to the nearest integer.  Double precision nominated candidates; the displayed residuals were recomputed with 100 decimal digits.  The table is exploratory only and plays no role in any proof.

\begin{center}
\begin{tabular}{r r c c}
\toprule
\(A\) & nearest \(B\) & \(\lvert A^\alpha-B\rvert\) &
\(\left\lvert\log_2A-\log_3B\right\rvert\)\\
\midrule
189,427 & 231,498,127 & $4.724862e-7$ & $1.857792e-15$ \\
378,854 & 694,494,381 & $1.417459e-6$ & $1.857792e-15$ \\
566,702 & 1,314,775,264 & $2.12226e-6$ & $1.469274e-15$ \\
306,982 & 497,586,056 & $2.213968e-6$ & $4.050034e-15$ \\
252,247 & 364,495,072 & $3.148417e-6$ & $7.862419e-15$ \\
958,460 & 3,023,921,230 & $4.140205e-6$ & $1.246255e-15$ \\
556,061 & 1,275,861,768 & $4.237204e-6$ & $3.022952e-15$ \\
757,708 & 2,083,483,143 & $4.252376e-6$ & $1.857792e-15$ \\
462,859 & 953,945,930 & $5.520422e-6$ & $5.267494e-15$ \\
613,964 & 1,492,758,168 & $6.641904e-6$ & $4.050034e-15$ \\
490,481 & 1,045,737,723 & $7.156446e-6$ & $6.22917e-15$ \\
596,949 & 1,427,722,864 & $7.260085e-6$ & $4.628639e-15$ \\
\bottomrule
\end{tabular}
\end{center}

Two lessons are robust.  First, absolute near misses of order roughly inverse to the search length are unsurprising; they do not signal a hidden exact relation.  Second, raising a near miss to powers amplifies rather than suppresses its logarithmic error.  A computational search should therefore be used to calibrate interpolation determinants or to locate special support patterns, not as evidence that a sufficiently close pair is ``almost'' a counterexample.

A more relevant future search would enumerate prime-exponent patterns in the canonical normal form, enforce content one and \(\min(a_2,b_3)=0\), split rank three from rank four, and use certified ball arithmetic for the star period.  This reduces duplicate work caused by integer shifts and perfect-power inflations.

\section{A concrete research program}

\subsection{Track A: prove a subquadratic polylogarithmic count}

The exact target is \(N_\alpha(H)=o((\log H)^2)\).  Potential tools include:
\begin{itemize}[leftmargin=2em]
\item a determinant method adapted to the multiplicative recurrence grid;
\item sharp rational-point bounds for the unrestricted exponential structure on the single curve \(Y=X^\alpha\);
\item a gap principle using two consecutive Hilbert-basis directions rather than isolated points;
\item a zero estimate that charges multiplicity in both semigroup directions.
\end{itemize}
The key advantage is logical: any exponent below two suffices.  One does not need the conjecturally optimal linear logarithmic count.

\subsection{Track B: degree-two Kummer orbit}

Formalize the Tannakian object generated by the Kummer extensions of the primes in \(S\) and isolate the matrix coefficient corresponding to \(F=X_2L_B-X_3L_A\).  The goal is not the entire period conjecture but the following implication for one external \(q\):
\[
 F(2\pi i,\log p)=0
 \quad\Longrightarrow\quad
 (\Omega\partial_qF)(2\pi i,\log p)=0.
\]
The right side contradicts unique factorization.  A proof may come from an orbit-dimension theorem, a degree-two period conjecture for a tiny mixed-Tate subquotient, or a comparison of matrix coefficients across realizations.

\subsection{Track C: isolate rank three}

In rank three, \cref{eq:mobius} gives
\[
 x=\frac{r+s\alpha}{u+v\alpha}
\]
with rational coefficients and a common external integer \(C\).  This case should be attacked separately.  Useful intermediate statements would be:
\begin{enumerate}[label=(\roman*)]
\item rule out the affine subcases \(u=0\) or \(v=0\);
\item prove algebraic independence of \(\alpha\) and any admissible nontrivial fractional-linear transform that simultaneously exponentiates at bases \(2\) and \(3\) to rationals;
\item exploit the primitive perfect-power condition on the common external part \(C\).
\end{enumerate}
Even a rank-three theorem would be meaningful progress, because it would force every counterexample into the generic rank-four case with two independent external divisor directions.

\subsection{Track D: formalize the unconditional core}

The semigroup theorem is well suited to Lean.  Almost all of it is finitely supported integer linear algebra once the six exponentials theorem and Gelfond--Schneider consequence are exposed as imported hypotheses.  A formalization would provide a reliable kernel for later experimental lemmas and prevent repeated mistakes about saturation, shifts, and root extraction.

\section{Lean formalization blueprint}

A practical dependency order is as follows.

\subsection{Algebraic layer}

Use finitely supported functions from primes to \(\Z\) for \(\calD\).  Define
\[
 \texttt{DivPair}:=\calD\times\calD,
\]
and the distinguished vector \(g_0=(e_2,e_3)\).  Prove:
\begin{itemize}[leftmargin=2em]
\item injectivity of the divisor map on positive rationals;
\item the gcd-of-minors formula for a two-column lattice;
\item the explicit content formula \eqref{eq:content};
\item root descent \cref{cor:root-descent};
\item equivalence between content one and modular independence for every \(\ell\).
\end{itemize}

\subsection{Analytic interface}

Define a structure carrying a real exponent \(t\), two positive rationals \(R,S\), and proofs
\[
 R=2^t,
 \qquad S=3^t.
\]
The only transcendence inputs needed for the structural theorem are:
\begin{enumerate}[label=(\alph*)]
\item \(\log2\) and \(\log3\) are \(\Q\)-linearly independent;
\item \(\alpha=\log3/\log2\) is transcendental;
\item the rank-cap consequence of six exponentials.
\end{enumerate}
These can initially be axiomatized in a namespace and later connected to formalized transcendence results.

\subsection{Semigroup layer}

Formalize the saturated lattice
\[
 \Lambda=(\Q g_0+\Q g_x)\cap F
\]
and select a basis extending \(g_0\).  Prove the sign lemmas coordinatewise:
\begin{itemize}[leftmargin=2em]
\item an external positive coordinate forces the coefficient of \(g_\xi\) to be nonnegative;
\item a zero distinguished coordinate forces the coefficient of \(g_0\) to be nonnegative.
\end{itemize}
The final theorem should state an equivalence between a nonintegral element and a monoid isomorphism
\[
 \N_0\times\N_0\cong\Eint,
 \qquad (m,n)\mapsto m+n\xi.
\]

\subsection{Tensor layer}

Implement the sparse tensor directly as a finitely supported symmetric pair map rather than importing a large symmetric-algebra hierarchy.  Prove:
\begin{itemize}[leftmargin=2em]
\item the coefficient expansion \eqref{eq:star-expanded};
\item formal-zero classification;
\item equality of tensor content and lattice index;
\item the external-prime derivative formula \eqref{eq:q-derivative}.
\end{itemize}
The final period-orbit theorem can then be represented as one explicit missing axiom.  This makes the exact frontier machine-checkable.

\section{Failure modes that the new normal form exposes}

\subsection{Replacing the real logarithm by local logarithms}
As explained above, there is no common local exponent.  Any proof doing this without a comparison theorem assumes its conclusion.

\subsection{Using only the determinant identity}
The control family has the same identity.  A valid proof must touch a nonzero coefficient of the primitive star tensor, preferably at an external prime.

\subsection{Choosing a least counterexample by size}
There need not be a least transcendental exponent under useful operations.  Saturation content, not numerical size, is the correct well-founded descent parameter.

\subsection{Ignoring integer-shift gauge}
The pairs \((A,B)\) and \((2^kA,3^kB)\) encode the same star tensor.  Arguments sensitive to \(a_2\) and \(b_3\) separately often prove only a statement about a chosen gauge.  The invariant is \(b_3-a_2\).

\subsection{Assuming an external prime gives a third base-exponent rectangle}
A prime dividing \(A\) or \(B\) does not imply that its \(x\)-th power is algebraic.  Thus six exponentials cannot be applied by simply adding that prime as a third base.  The external prime is useful through valuations, saturation, or Kummer directions, not through an unproved algebraicity of \(q^x\).

\subsection{Treating numerical proximity as algebraic evidence}
The relevant residual is a computable real number but can be very small for accidental reasons.  Without a proven lower bound for the particular quadratic logarithmic form, finite precision is not a certificate of nonvanishing over an unbounded family.

\section{Conclusions}

The central outcome is a cleaner frontier.

If the conjecture is false, there is a unique canonical irrational \(\xi\) for which the entire integer exponent set is the free monoid
\[
 \Eint=\N_0+\N_0\xi.
\]
The joint divisor vectors of \(1\) and \(\xi\) form a saturated rank-two lattice.  Their Pl\"ucker content is one, one distinguished base valuation vanishes, and an external prime is unavoidable.  The corresponding quadratic relation is a nonzero primitive two-star tensor in the kernel of the product-of-logarithms period map.

This yields two focused ways forward:
\begin{enumerate}[label=(\arabic*)]
\item rule out the quadratic logarithmic growth forced by the free semigroup, i.e. prove
\(N_\alpha(H)=o((\log H)^2)\);
\item prove that the zero two-star period must be stable under one external Kummer translation, whose derivative is visibly nonzero.
\end{enumerate}
Both targets exploit information absent from the bare four-exponentials determinant: the first uses the full recurrence grid, and the second uses the external-prime support of the primitive tensor.

The present report does not close the final transcendence gap.  It does, however, turn an amorphous search into a canonical normal form with a small number of exact interfaces.  Any future claimed proof can now be stress-tested against saturation, gauge invariance, the external-prime derivative, the rank-three normal form, and the quadratic counting lower bound before substantial formalization effort is invested.

\appendix

\section{A compact equivalence dictionary}

For quick reference, the following data encode the same hypothetical counterexample.

\begin{center}
\begin{longtable}{>{\raggedright\arraybackslash}p{0.23\textwidth} p{0.69\textwidth}}
\toprule
Language & Object\\
\midrule
Exponent & An irrational \(\xi\) with \(2^\xi,3^\xi\in\N\).\\
Integer point & \((C,D)=(2^\xi,3^\xi)\) on \(D=C^\alpha\).\\
Log matrix & A rank-one real matrix \(\bigl(\begin{smallmatrix}\log2&\log C\\\log3&\log D\end{smallmatrix}\bigr)\) with both rational row and column ranks two.\\
Divisor lattice & A saturated rank-two lattice generated by \((e_2,e_3)\) and \((\divv C,\divv D)\).\\
Affine semigroup & \(\Eint\cong\N_0^2\) via \((m,n)\mapsto m+n\xi\).\\
Tensor period & A primitive nonzero two-star tensor \(e_2\tensorstar\divv D-e_3\tensorstar\divv C\) with zero real period.\\
Counting & \(N_\alpha(H)\sim(2\xi(\log3)^2)^{-1}(\log H)^2\).\\
Kummer witness & For every \(\ell\), the two divisor vectors are independent modulo \(\ell\); at an external \(q\), the orbit derivative is \(b_q\log2-a_q\log3\ne0\).\\
\bottomrule
\end{longtable}
\end{center}

\section{Derivation of the content formula by coordinates}

Order the basis of \(F=\calD\oplus\calD\) as \(e_{1,p}\) and \(e_{2,p}\).  The two columns are
\[
 g_0=e_{1,2}+e_{2,3},
 \qquad
 g_x=\sum_pa_pe_{1,p}+\sum_pb_pe_{2,p}.
\]
The wedge is
\[
 g_0\wedge g_x
 =e_{1,2}\wedge g_x+e_{2,3}\wedge g_x.
\]
Its coordinates include \(a_p\) for \(p\ne2\), \(b_p\) for \(p\ne3\), and at the basis pair \((e_{1,2},e_{2,3})\) the coefficient is \(b_3-a_2\).  All remaining nonzero coordinates repeat these numbers up to sign.  Their gcd is therefore \(\delta(A,B)\).  The Smith normal form of a two-column integer matrix identifies this gcd with the index in the saturated plane.

The symmetric tensor \(\Sigma(A,B)\) has exactly the same irredundant coefficient list.  This is why lattice saturation, root extraction, and tensor primitivity are the same operation viewed in three languages.

\section{The control family and antisymmetry}

Before symmetrization, the raw tensor associated with the determinant is
\[
 T=e_2\otimes b-e_3\otimes a.
\]
For the control pair \(a=me_2\), \(b=me_3\),
\[
 T=m(e_2\otimes e_3-e_3\otimes e_2),
\]
which is nonzero but purely antisymmetric.  Multiplication of real numbers kills every antisymmetric tensor automatically.  Passing to \(\Sym^2(V)\) quotients out this universal commutativity kernel and makes the control tensor literally zero.  A counterexample is therefore not merely another element of the obvious kernel; it creates a new symmetric quadratic relation.

This observation is useful when designing period objects.  Working in the full tensor square without first removing the exterior square allows the control family to masquerade as the same obstruction one is trying to exclude.

\section{References}

\begin{thebibliography}{99}

\bibitem{AE1944}
L. Alaoglu and P. Erd\H{o}s,
\emph{On highly composite and similar numbers},
Transactions of the American Mathematical Society \textbf{56} (1944), 448--469.

\bibitem{Baker}
A. Baker,
\emph{Transcendental Number Theory},
Cambridge University Press, 1975.

\bibitem{Lang}
S. Lang,
\emph{Introduction to Transcendental Numbers},
Addison--Wesley, 1966.

\bibitem{PW}
J. Pila and A. J. Wilkie,
\emph{The rational points of a definable set},
Duke Mathematical Journal \textbf{133} (2006), 591--616.

\bibitem{Waldschmidt}
M. Waldschmidt,
\emph{Diophantine Approximation on Linear Algebraic Groups: Transcendence Properties of the Exponential Function in Several Variables},
Springer, 2000.

\bibitem{Wustholz}
G. W\"ustholz,
\emph{Algebraische Punkte auf analytischen Untergruppen algebraischer Gruppen},
Annals of Mathematics \textbf{129} (1989), 501--517.

\bibitem{KZ}
M. Kontsevich and D. Zagier,
\emph{Periods}, in \emph{Mathematics Unlimited---2001 and Beyond},
Springer, 2001, 771--808.

\end{thebibliography}

\end{document}
